\begin{ouverture}[conditions]{13}{Algorithmique, logique et raisonnement}{Outils transversaux}{Fonctions Python · Listes · Ensembles · Preuve}
\objectif{Décomposer une tâche en fonctions testables.}
\objectif{Utiliser variables, boucles, conditions et listes.}
\objectif{Manipuler ensembles et propositions.}
\objectif{Distinguer implication, réciproque et équivalence.}
\prerequis{Algorithmique de Seconde et Première ; notions rencontrées dans les chapitres.}
\end{ouverture}
\essentiel

\essSection{1}{Écrire un programme qui répond à une question}
Une fonction informatique possède des paramètres, un traitement et une valeur renvoyée par \code{return}. Définir ce qui entre et ce qui sort aide à vérifier le programme.
L'affectation \code{u = 2*u + 1} remplace la valeur de \code{u} ; ce n'est pas une équation à résoudre.
Une boucle \code{for} convient à un nombre de répétitions connu ; une boucle \code{while} dépend d'une condition dont il faut justifier qu'elle finit par devenir fausse.
\begin{python}
def termes_affines(u0, a, b, n):
    termes = [u0]
    u = u0
    for k in range(n):
        u = a*u + b
        termes.append(u)
    return termes
\end{python}
Cette fonction rend $n+1$ valeurs, de $u_0$ à $u_n$. Les indices d'une liste Python commencent à zéro.
La programmation modulaire évite de réécrire les mêmes calculs : une fonction pour une masse binomiale, une autre pour un cumul, une troisième pour chercher un seuil.
\essSection{2}{Expérimenter, puis vérifier}
Pour un calcul numérique, tester un cas connu, un cas limite et les hypothèses de domaine. Des nombres décimaux sont représentés approximativement : une égalité mathématique n'est pas toujours une égalité machine exacte.
\begin{savoirfaire}[Un compte rendu d'algorithme]
Préciser les entrées ; expliquer chaque variable ; donner la condition d'arrêt ; tester sur un exemple dont on connaît le résultat ; interpréter la sortie et sa précision.
Pour une simulation, préciser le modèle et le nombre de répétitions. Pour un encadrement, vérifier que la racine reste dans l'intervalle.
\end{savoirfaire}
\figurelivre{c13-algorithme.pdf}{La boucle d'un algorithme de seuil : initialisation, test, mise à jour et sortie. Une condition d'arrêt doit être cohérente avec l'inégalité demandée.}
\essSection{3}{Ensembles et événements}
$x\in A$ signifie que $x$ est un élément de $A$ ; $A\subset B$ que chaque élément de $A$ appartient à $B$.
$A\cap B$ signifie « dans $A$ et dans $B$ » ; $A\cup B$, « dans $A$ ou dans $B$, éventuellement les deux ».
Le complémentaire de $A$ dans $E$ est $E\setminus A$, ou $\bar A$ en probabilités.
Un couple $(x,y)$ est ordonné : en général $(x,y)\ne(y,x)$.
\begin{exemple}
Pour $A=[0,3]$ et $B=[2,5]$, $A\cap B=[2,3]$ et $A\cup B=[0,5]$.
La négation de « $x>0$ et $x<3$ » est « $x\leq0$ ou $x\geq3$ ».
\end{exemple}
\essSection{4}{Choisir une forme de raisonnement}
Une proposition est un énoncé auquel on peut attribuer vrai ou faux lorsque ses variables et son domaine sont précisés.
L'implication « si $P$, alors $Q$ » ne dit pas que $Q$ entraîne $P$. Sa réciproque est « si $Q$, alors $P$ ».
L'équivalence nécessite les deux implications.
Un contre-exemple suffit à réfuter une affirmation universelle ; des exemples, même nombreux, ne suffisent pas à la prouver.
\begin{exemple}[Une réciproque fausse]
Si $x=2$, alors $x^2=4$. La réciproque est fausse sur $\R$ : $x=-2$ satisfait $x^2=4$.
En revanche, sur $[0,+\infty[$, $x^2=4$ équivaut à $x=2$.
\end{exemple}
\essSection{5}{Quantifier et communiquer}
« Pour tout $x>0$, $\ln x\leq x-1$ » est universel. « Il existe $x>0$ tel que $\ln x=0$ » est existentiel, vérifié par $x=1$.
Le symbole $\sum$ abrège une addition : $\sum_{k=0}^{n}u_k$ contient $n+1$ termes.
À l'oral, annoncer l'objectif, les hypothèses, la méthode, le résultat et sa limite de validité. Un graphique doit présenter ses axes et une légende ; un nombre doit garder son unité ou son statut de probabilité.

\exercices
\markright{Prendre confiance}
\begingroup\renewcommand{\otBesoinPlace}[1]{}\begin{minipage}{\linewidth}
\rubrique{Prendre confiance}
\begin{exercice}[1]{Affectation}
On exécute \code{u = 3}, puis \code{u = 2*u + 1}.
\begin{questions}
\item Quelle est la valeur finale ?
\item Pourquoi ne résout-on pas $u=2u+1$ ?
\item Que donnerait une répétition supplémentaire ?
\end{questions}
\end{exercice}
\begin{exercice}[1]{Liste}
On appelle \code{termes_affines(0, 2, 1, 3)}.
\begin{questions}
\item Donner la liste rendue.
\item Quelle est sa longueur ?
\item Quelle valeur occupe l’indice $2$ ?
\end{questions}
\end{exercice}
\begin{exercice}[1]{Ensembles}
$A=\{1,2,3\}$ et $B=\{3,4\}$ dans $E=\{1,2,3,4,5\}$.
\begin{questions}
\item Donner $A\cap B$ et $A\cup B$.
\item Donner $E\setminus A$.
\item L’affirmation $A\subset B$ est-elle vraie ?
\end{questions}
\end{exercice}
\begin{exercice}[1]{Nier une proposition}

\begin{questions}
\item Nier « $x>2$ ».
\item Nier « $x<0$ ou $x>4$ ».
\item Nier « $x\geq1$ et $x\leq3$ ».
\end{questions}
\end{exercice}
\end{minipage}\endgroup
\clearpage
\markright{Appliquer et relier}
\begingroup\renewcommand{\otBesoinPlace}[1]{}\begin{minipage}{\linewidth}
\rubrique{Appliquer et relier}
\begin{exercice}[2]{Finir une boucle}
On part de $u=1$ et l’on multiplie par $2$ tant que $u<100$.
\begin{questions}
\item Combien de mises à jour sont faites ?
\item Quelle est la valeur finale ?
\item Que changerait la condition $u\leq128$ ?
\end{questions}
\end{exercice}
\begin{exercice}[2]{Une fonction de moyenne}

\begin{questions}
\item Écrire une fonction Python donnant la moyenne d’une liste non vide.
\item Quelle hypothèse faut-il vérifier ?
\item Tester sur $[2,4,9]$.
\end{questions}
\end{exercice}
\begin{exercice}[2]{Implication}

\begin{questions}
\item La proposition « $x>2$ entraîne $x^2>4$ » est-elle vraie sur $\R$ ?
\item Écrire la réciproque et la tester.
\item Quel domaine rend les deux propositions équivalentes ?
\end{questions}
\end{exercice}
\begin{exercice}[2]{Tester un programme de cumul}
Une fonction est censée calculer $P(X\leq k)$ pour une binomiale de taille $n$.
\begin{questions}
\item Que doit-elle rendre pour $k=n$ ?
\item Que doit-elle rendre pour $k=0$ ?
\item Pourquoi un résultat légèrement différent de $1$ en machine peut-il apparaître ?
\end{questions}
\end{exercice}
\end{minipage}\endgroup
\clearpage
\markright{Chercher et argumenter}
\begingroup\renewcommand{\otBesoinPlace}[1]{}\begin{minipage}{\linewidth}
\rubrique{Chercher et argumenter}
\begin{exercice}[3]{Contre-exemple}

\begin{questions}
\item Réfuter « toute suite positive est croissante ».
\item Réfuter « si $f\,\!^{\prime}(a)=0$, alors $f$ a un extremum en $a$ ».
\item Pourquoi un seul exemple suffit-il dans chaque cas ?
\end{questions}
\end{exercice}
\begin{exercice}[3]{Raisonner sur un intervalle}
On cherche les solutions de $\ln(x-1)=\ln(3-x)$.
\begin{questions}
\item Déterminer le domaine commun.
\item Résoudre par équivalence.
\item Comment présenter une preuve complète ?
\end{questions}
\end{exercice}
\begin{exercice}[3]{Somme : une erreur de borne}
Un programme utilise \code{sum(k for k in range(n))}.
\begin{questions}
\item Quelle somme calcule-t-il ?
\item Comment calculer $1+\cdots+n$ ?
\item Comparer pour $n=5$.
\end{questions}
\end{exercice}
\begin{exercice}[3]{Un oral argumenté}
Préparer une présentation de trois minutes sur la question « Un test très sensible donne-t-il toujours un résultat positif fiable ? ».
\begin{questions}
\item Choisir un modèle chiffré.
\item Identifier le schéma et le calcul à présenter.
\item Formuler une conclusion et une limite.
\end{questions}
\end{exercice}
\end{minipage}\endgroup
